\documentclass[10pt]{article}
\usepackage[a4paper,margin=24mm]{geometry}
\usepackage{amsmath,amssymb,amsthm,microtype,hyperref}
\hypersetup{colorlinks=true,urlcolor=blue,linkcolor=black}
\newtheorem{theorem}{Theorem}
\newtheorem{lemma}[theorem]{Lemma}
\newtheorem{corollary}[theorem]{Corollary}
\newcommand{\Var}{\operatorname{Var}}
\newcommand{\IS}{\mathcal I}
\newcommand{\Z}{Z}
\title{The entropy cost of a linear log-concavity failure run}
\author{Research note prepared for Jan Mikul\'ik with ChatGPT}
\date{30 September 2026}
\begin{document}
\maketitle
\begin{abstract}
The degree-free forest obstruction extends to graphs whose cyclic part has
subpolynomially many independent configurations. Define
\(\kappa(G)=\min_{G-S\text{ a forest}}\log_2 Z_{G[S]}(1)\).
If \(\kappa(G_n)=o(\log n)\), every consecutive strict log-concavity
failure run in \(Z_{G_n}\) has length \(o(n)\). In particular, this holds
when the feedback vertex number is \(o(\log n)\), and when deleting a clique
of size \(n^{o(1)}\) leaves a forest. Edges incident with the deleted part
are unrestricted. The proof uses an exact marginal reweighting identity,
not an assumption about the conditional independence of cyclic vertices.
\end{abstract}

\noindent\textbf{Scope.} This is a proved extension of the supplied forest
argument, rederived below with an optimized variance exponent. It is a
research draft: independent review and publication priority are not
established. The accompanying exact finite audit checks the new transfer
identities; enumeration is not a substitute for the universal proof.

\section{Statement}
All graphs are finite and simple. Write
\[
 Z_G(\lambda)=\sum_{I\in\IS(G)}\lambda^{|I|}
   =\sum_{j=0}^{\alpha(G)}i_j(G)\lambda^j.
\]
A strict failure at index \(j\) means
\(i_j^2<i_{j-1}i_{j+1}\). A run at \(u,\ldots,v\) has length
\(L=v-u+1\); put \(d=L+1\). Set
\[
 \kappa(G)=\min_{S\subseteq V(G):\,G-S\text{ is a forest}}
               \log_2 Z_{G[S]}(1).
\]
The empty graph counts as a forest, and \(Z_\varnothing=1\).

\begin{theorem}[An explicit structural obstruction]\label{thm:main}
Fix \(0<\delta\le1\), and define
\[
 \Lambda_\delta=4^{1/\delta},\qquad
 B_\delta=2\sqrt{\Lambda_\delta},\qquad
 \eta_\delta=
 \frac{\log(1+1/\Lambda_\delta)}
      {\log B_\delta+\log(1+1/\Lambda_\delta)} >0.
\]
If an \(n\)-vertex graph has a strict failure run with \(L+1\ge\delta n\),
then
\begin{equation}\label{eq:barrier}
\boxed{\quad
 \kappa(G)>
 \min\left\{\frac{\delta n}{\sqrt{12}},
 \frac{\eta_\delta\log n-\log(36/\delta^2)}
      {\log(2\Lambda_\delta)}\right\}.
 \quad}
\end{equation}
Here and below unmarked logarithms are natural. In particular, for every
fixed \(\delta\), a run of this length forces
\(\kappa(G)\ge c_\delta\log n\) for all sufficiently large \(n\), where
\(c_\delta=\eta_\delta/[2\log(2\Lambda_\delta)]\).
\end{theorem}

\begin{corollary}\label{cor:asym}
For any sequence of graphs of orders \(n\to\infty\),
\[
 \kappa(G_n)=o(\log n)
 \quad\Longrightarrow\quad
 \frac{\text{longest strict failure run in }G_n}{n}\longrightarrow0.
\]
There are no maximum-degree, edge-count, or component-count assumptions.
\end{corollary}

\section{The exact transfer identity}
Allow positive vertex activities \(\lambda_v\). For any graph \(H\),
\(Z_H=\sum_{I\in\IS(H)}\prod_{v\in I}\lambda_v\).
Fix any \(S\subseteq V(G)\), let \(F=G-S\), and let \(P_F\) denote the
hard-core measure on \(F\). For \(J\in\IS(F)\), define
\[
 W(J)=Z_{G[S\setminus N_G(J)]},\qquad C=Z_{G[S]}.
\]
Every independent set in \(G\) is uniquely \(J\cup K\), where
\(K\in\IS(G[S\setminus N_G(J)])\). Consequently
\begin{equation}\label{eq:density}
 P_G(I\cap V(F)=J)=\frac{P_F(J)W(J)}{\mathbb E_F W},
 \qquad 1\le W(J)\le C.
\end{equation}
This is an exact identity. No acyclicity is needed here.

\begin{lemma}[Variance transfer]\label{lem:transfer}
Let \(X_G=|I|\), \(X_F=|J|\), and \(a=\alpha(G[S])\). Then
\begin{equation}\label{eq:transfer}
 \sqrt{\Var_G X_G}\le
       \sqrt{C\,\Var_F X_F}+\frac a2.
\end{equation}
More generally, every nonnegative observable \(f(J)\) satisfies
\(\mathbb E_G f(I\cap F)\le C\mathbb E_F f(J)\).
\end{lemma}
\begin{proof}
The expectation inequality follows from \eqref{eq:density}, since
\(\mathbb E_F W\ge1\). Under \(P_G\), write \(X_G=Y+T\), with
\(Y=|I\cap F|\) and \(0\le T=|I\cap S|\le a\). Choosing
\(b=\mathbb E_F X_F\) gives
\[
 \Var_GY=\min_{t\in\mathbb R}\mathbb E_G(Y-t)^2
 \le\mathbb E_G(Y-b)^2\le C\Var_F X_F.
\]
Also \(\Var_GT\le a^2/4\): a variable in \([0,a]\) has variance at most
\(a^2/4\). The triangle inequality in \(L^2(P_G)\) for the centered
variables \(Y\) and \(T\) proves \eqref{eq:transfer}. These variables
need not be independent.
\end{proof}

\section{A self-contained forest variance estimate}
\begin{lemma}\label{lem:forest}
Let \(F\) be a forest with at most \(n\ge1\) vertices and positive
activities at most \(\Lambda\ge1\). Put
\[
 B=2\sqrt\Lambda,\quad q=\frac{\Lambda}{1+\Lambda},\quad
 \eta=\frac{\log(1/q)}{\log(B/q)}.
\]
Then, without a lower threshold on \(n\),
\begin{equation}\label{eq:variance}
 \Var_F X_F\le\frac34 n^{2-\eta}.
\end{equation}
\end{lemma}
\begin{proof}
The empty forest is immediate. Otherwise let \(S_v\) indicate occupancy,
\(\mu_v=\mathbb E S_v\), \(s_v=\sqrt{\Var S_v}\), and
\(w_v=\mu_v/(1-\mu_v)\). For an oriented edge \(u\to v\), let
\(R_{u\to v}\) be the occupancy odds at \(u\) in the component obtained
by deleting \(uv\). Conditioning at its root gives
\[
 R_{u\to v}=\lambda_u\prod_{z\sim u,\,z\ne v}(1+R_{z\to u})^{-1},
 \qquad w_v=\lambda_v\prod_{u\sim v}(1+R_{u\to v})^{-1}.
\]
In particular \(R_{u\to v}\le\Lambda\). Adjacent vertices cannot both
be occupied, so, with \(a_{uv}=|\operatorname{Corr}(S_u,S_v)|\),
\[
 a_{uv}^2=w_uw_v.
\]
Fix \(v\) and write \(x_u=R_{u\to v}\). Edge conditioning gives
\(w_u=x_u/(1+R_{v\to u})\le x_u\). Therefore
\[
 \sum_{u\sim v}a_{uv}^2
 \le\lambda_v\frac{\sum_{u\sim v}x_u}{\prod_{u\sim v}(1+x_u)}
 \le\lambda_v\le\Lambda.
\]
The middle product is at least \(1+\sum x_u\). For an edge, writing
\(R=R_{u\to v}\) and \(T=R_{v\to u}\) also gives
\[
 a_{uv}^2=\frac{R}{1+R}\frac{T}{1+T}\le q^2.
\]
Let \(A\) be weighted adjacency with entries \(a_{uv}\). Root each
component, and let \(D\) contain the parent-to-child entries of \(A\).
Distinct rows of \(D\) have disjoint supports. Thus \(DD^\mathsf T\)
is diagonal, with entries at most \(\Lambda\), and
\(\|A\|_2=\|D+D^\mathsf T\|_2\le2\sqrt\Lambda=B\).

For vertices at distance \(r\) in the same component, the tree Markov
property implies
\[
 \operatorname{Cov}(S_u,S_v)=(-1)^r s_us_v
                  \prod_{e\in P(u,v)}a_e.
\]
For completeness, a centered function of a binary variable is a scalar
multiple of its centered indicator. Hence on an edge the conditional
expectation of the normalized centered indicator is its correlation
times the normalized centered indicator at the other endpoint.
Composing this relation along the path proves the formula. Distinct
components are independent.

The simple path is one of the nonnegative walks counted in \(A^r\).
Writing \(y_v=s_v\), the sum over ordered pairs at distance \(r\) is
at most \(y^\mathsf T A^r y\le(n/4)B^r\). Each individual pair at
distance \(r\) has absolute covariance at most \(q^r/4\). Therefore,
for every integer \(R\ge0\),
\begin{equation}\label{eq:radius}
 \Var_F X_F\le \frac n4\sum_{r=0}^R B^r+
                      \frac{n^2}4q^{R+1}
 \le\frac n2 B^R+\frac{n^2}4q^{R+1}.
\end{equation}
The exponent \(R+1\) uses the fact that all remaining distances exceed
\(R\). Put \(t=\log n/\log(B/q)\) and \(R=\lfloor t\rfloor\).
Then
\[
 B^R\le B^t=n^{1-\eta},\qquad q^{R+1}\le q^t=n^{-\eta}.
\]
Substitution in \eqref{eq:radius} proves \eqref{eq:variance}.
\end{proof}

\section{A failure window and its centered activity}
\begin{lemma}\label{lem:window}
For any \(n\)-vertex graph with a strict failure run at \(u,\ldots,v\),
put \(c=(u+v)/2\) and \(d=v-u+2\). There is a unique common activity
\(\lambda>0\) with \(\mathbb E_\lambda X_G=c\), and at that activity
\begin{equation}\label{eq:window}
 \Var_\lambda X_G>\frac{d^2}{12},\qquad
 \lambda\le 2^{2n/d}.
\end{equation}
\end{lemma}
\begin{proof}
Every coefficient from \(0\) through \(\alpha(G)\) is positive, since
subsets of a maximum independent set are independent. The mean tends
to \(0\) and \(\alpha(G)\) as \(\lambda\) tends to \(0\) and
\(\infty\), respectively. Its derivative with respect to \(\log\lambda\)
is the positive variance. Also \(0<c<\alpha(G)\).

At the centered activity let \(p_j=i_j\lambda^j/Z_G(\lambda)\).
Tilting preserves the strict failure inequalities. The arithmetic-geometric
mean inequality implies \(p_{j-1}-2p_j+p_{j+1}>0\) for
\(u\le j\le v\). For any numbers \(z_0,\ldots,z_d\), collecting
coefficients gives the exact identity
\begin{align*}
 \sum_{j=0}^{d}\left[(j-d/2)^2-\frac{d(d+2)}{12}\right]z_j
 =\sum_{j=1}^{d-1}
 \frac{j(j+1)(d-j)(d-j+1)}{12}
                  (z_{j-1}-2z_j+z_{j+1}).
\end{align*}
Take \(z_j=p_{u-1+j}\). The right side is nonnegative, and every index
outside the window has squared distance from \(c\) larger than
\(d(d+2)/12\). Summing all probabilities gives
\(\Var X_G\ge d(d+2)/12>d^2/12\).

The Shannon entropy of the independent-set distribution, in nats, is
\(H=\log Z_G(\lambda)-c\log\lambda\le n\log2\).
Since \(Z_G(\lambda)\ge\lambda^{\alpha(G)}\),
\((\alpha(G)-c)\log\lambda\le n\log2\).
The enlarged window ends at \(v+1\le\alpha(G)\), so
\(\alpha(G)-c\ge d/2\). If \(\log\lambda\ge0\) this yields the stated
bound; if \(\lambda<1\), the bound is automatic.
\end{proof}

\section{Proof of the entropy barrier}
\begin{proof}[Proof of Theorem \ref{thm:main}]
Choose any \(S\) with \(F=G-S\) a forest, and put
\(h=\log_2 Z_{G[S]}(1)\) and \(a=\alpha(G[S])\).
All subsets of a maximum independent set in \(G[S]\) contribute to
\(Z_{G[S]}(1)\), so \(a\le h\). At the centered activity of
Lemma \ref{lem:window}, \(\lambda\le\Lambda_\delta\). Consequently
\begin{equation}\label{eq:corecost}
 C=Z_{G[S]}(\lambda)\le Z_{G[S]}(1)\Lambda_\delta^a
                   \le(2\Lambda_\delta)^h.
\end{equation}
Apply Lemmas \ref{lem:transfer} and \ref{lem:forest}, using the ambient
order \(n\) as an upper bound on the forest order. We obtain
\begin{equation}\label{eq:master}
 \frac{\delta}{\sqrt{12}}
 <\frac{\sqrt{\Var_GX_G}}n
 \le\frac{\sqrt3}{2}(2\Lambda_\delta)^{h/2}n^{-\eta_\delta/2}
                         +\frac{h}{2n}.
\end{equation}
If \(h>\delta n/\sqrt{12}\), the desired minimum bound follows.
Otherwise the last term is at most \(\delta/(2\sqrt{12})\).
Rearranging and squaring the remaining positive terms gives
\[
 (2\Lambda_\delta)^h>\frac{\delta^2}{36}n^{\eta_\delta}.
\]
Taking logarithms proves the other alternative in \eqref{eq:barrier}.
The bound holds for every admissible \(S\), including one attaining
the minimum defining \(\kappa(G)\). For fixed \(\delta\), the logarithmic
term is eventually the smaller term and eventually at least
\(c_\delta\log n\).
\end{proof}

\begin{proof}[Proof of Corollary \ref{cor:asym}]
If the conclusion failed, some fixed \(\delta>0\) and an infinite
subsequence would have \(L+1\ge\delta n\). Theorem \ref{thm:main} would
force \(\kappa(G_n)\ge c_\delta\log n\) along that subsequence,
contradicting \(\kappa(G_n)=o(\log n)\).
\end{proof}

\section{Concrete consequences}
\begin{corollary}[Feedback vertices and a finite threshold]
Suppose deleting \(k\) vertices makes \(G\) a forest. Fix \(\delta\) and
use the parameters of Theorem \ref{thm:main}. A run with
\(L+1\ge\delta n\) is impossible whenever
\[
 n\ge\max\left\{1,\frac{\sqrt{12}\,k}{\delta},
       \left[\frac{36(1+\Lambda_\delta)^k}{\delta^2}
                  \right]^{1/\eta_\delta}\right\}.
\]
Thus feedback vertex number \(o(\log n)\) suffices for the asymptotic
obstruction.
\end{corollary}
\begin{proof}
For the chosen \(S\), \(a\le k\) and
\(C\le(1+\Lambda_\delta)^k\). Under the displayed threshold each of
the two upper terms in the analogue of \eqref{eq:master} is at most
\(\delta/(2\sqrt{12})\), contradicting its strict lower bound.
Also \(Z_{G[S]}(1)\le2^k\), so \(\kappa(G)\le k\).
\end{proof}

\noindent\textbf{Large dense cores.} If \(S\) induces a clique of size
\(s\), then \(Z_{G[S]}(1)=1+s\). Thus a clique of size \(s=n^{o(1)}\)
may be deleted, with every edge between the clique and the remaining
forest arbitrary, while the conclusion still holds. For example,
\(s=\lceil\exp(\sqrt{\log n})\rceil\) is far larger than \(\log n\),
but its configuration cost is only \(O(\sqrt{\log n})\).
More generally, \(Z_{G[S]}(1)=n^{o(1)}\) is the sufficient condition.

\medskip
\noindent\textbf{What has changed from the supplied proof.} The supplied
note treats forests. The new structural statement treats arbitrary
attachments to an induced cyclic core, charging that core by the number
of its independent configurations. The forest exponent is also optimized:
balancing \(B^R\) against \(q^{R+1}\) yields
\(\eta=\log(1/q)/\log(B/q)\) for every ambient order \(n\ge1\).
Neither the exact marginal identity nor the graph entropy barrier relies
on the finite audit.

\medskip
\noindent\textbf{Verification.} Run
\texttt{python verify\_core\_entropy\_barrier.py}.
It uses the standard library and exact rational arithmetic. It independently
enumerates graph independent sets, verifies the reweighting identity,
observable domination, the sharp square-root variance inequality without
floating-point roots, and the core counting inequalities. The JSON report
states the finite scope and the number of checks.

\medskip
\noindent\textbf{Literature boundary.} Spectral estimates on trees below
\(e^2\) appear in X.~Chen, X.~Yang, Y.~Yin and X.~Zhang,
\emph{Spectral Independence Beyond Total Influence on Trees and Related
Graphs}, \href{https://arxiv.org/abs/2404.04668}{arXiv:2404.04668}.
The supplied \emph{A degree-free obstruction to linear runs of
log-concavity failures in forests} (28 September 2026) is the starting
point for this note. No claim that the present extension is absent from
all existing literature is made.
\end{document}
